\chapter{Ramas de entrenamiento y de servicios}
\label{ap:ramas}

El programa que se instala en la llamada está en la rama principal y se describe en la Sección~\ref{sec:estructura_codigo}. Este apéndice reúne las otras ramas. El detalle de directorios, módulos y convenciones de nombres queda en el \texttt{README} de cada una.

\section{Repositorio de entrenamiento del clasificador}
Este repositorio entrena y evalúa el clasificador de señas aisladas. Vive en la rama \texttt{scratch-mediapipe-v3}. El archivo de pesos que sale de ese entrenamiento es el que después carga el motor local.

Todo el código está en \texttt{src/}. \texttt{train.py} ajusta el modelo y guarda el peso, \texttt{tune\_optuna.py} busca hiperparámetros, y \texttt{preprocessing.py} convierte los videos en coordenadas. \texttt{camera.py} mide el acierto frente a la cámara. La arquitectura está en \texttt{model\_arch.py}. \texttt{config.py} fija las rutas, las 100 clases y la ventana de 16 fotogramas. La extracción usa el Holistic del navegador, el mismo paquete que la extensión, a través de \texttt{holistic\_web.py}. Los videos quedan en \texttt{dataset/} y las coordenadas en \texttt{dataset\_landmarks\_16frames/}. Los pesos y los registros salen a \texttt{src/model\_new\_2026\_10\_20/}, con el checkpoint \texttt{tinyskeleton\_best.pth}.

\section{Servicio semántico en la nube}
Este repositorio es el servicio que el motor local consulta para traducir. Atiende las dos direcciones, de glosas hacia español y de español hacia glosas, con los binarios que entregan los repositorios de entrenamiento. Su tarea es responder un enunciado ya armado, y no forma parte de lo que se instala en la computadora de quien llama.

El arranque está en \texttt{run\_semantic\_server.py}. El servicio vive en \texttt{src/semantic/}: \texttt{api.py} expone las rutas, \texttt{service.py} separa los dos canales y \texttt{engine.py} ejecuta el binario. Los pesos no comparten carpeta. \texttt{models/sign/} guarda el canal de glosas hacia español y \texttt{models/hearing/} el canal inverso. En cada una hay un binario principal, un respaldo y el \texttt{sys\_prompt.txt} de esa dirección.

\section{Entrenamiento del canal directo}
Este repositorio ajusta y evalúa los modelos que pasan de glosas a una oración en español. El binario que resulta de ese entrenamiento es el que después carga el servicio en la nube para el canal directo.

Todo el código está en \texttt{src/semantic/}. \texttt{train.py} ajusta un modelo y exporta el binario, \texttt{eval\_trial.py} lo compara sobre oraciones fuera del corte, y \texttt{models.py} es el catálogo: sumar una familia es agregar una entrada ahí. El corpus está en \texttt{dataset\_glosses.json}, las oraciones de prueba en \texttt{dataset\_trial.json} y las instrucciones del modelo en \texttt{prompts/sys\_prompt.txt}. Los pesos y los registros salen a \texttt{outputs/}. Cada familia se identifica con un id estable, como \texttt{qwen2.5-1.5b} o \texttt{llama-3.2-1b}, y ese mismo id nombra su carpeta de salida.

\section{Entrenamiento del canal inverso}
Este repositorio ajusta y evalúa los modelos que pasan de una oración en español a una secuencia de glosas. El corpus y el binario son propios de esta dirección, y el servicio en la nube los usa para el canal inverso.

La disposición repite la del canal directo, también bajo \texttt{src/semantic/} y también con los módulos en inglés: \texttt{train.py} entrena, \texttt{eval\_trial.py} compara y \texttt{config.py} reúne rutas, hiperparámetros y la lista de modelos. El corpus, el archivo de prueba y el prompt ocupan los mismos nombres de archivo que en el otro laboratorio. Los id de modelo también coinciden, de modo que \texttt{qwen2.5-1.5b} designa la misma familia en las dos direcciones, con pesos distintos.

\section{Servicio de poses del avatar}
Este repositorio publica los archivos de pose que el motor local pide para dibujar el avatar. No entrena modelos ni traduce texto. El cliente de la videollamada solicita esas poses a este servicio.

El código cabe en \texttt{app/}: \texttt{main.py} arranca el servicio, \texttt{config.py} fija rutas y tipos de archivo, y \texttt{fingerspell.py} resuelve el alfabeto manual. Las poses están en \texttt{poses/}. El nombre del archivo, sin extensión, es la glosa ya normalizada: mayúsculas, sin tildes y en ASCII. Dentro de una palabra la eñe pasa a ene, y la letra eñe, cuando es la glosa entera, se pide como \texttt{ENIE}.
